% !TITLE 又呈吴郎
% !AUTHOR 杜甫
% !DYNASTY 唐
% !GENRE 七言律诗
% !ORDER 35

\begin{poem}
堂前扑枣\textsuperscript{1}任西邻，无食无儿一妇人。
不为困穷宁有此\textsuperscript{2}？只缘恐惧转须亲\textsuperscript{3}。
即防远客虽多事\textsuperscript{4}，便插疏篱却甚真\textsuperscript{5}。
已诉征求贫到骨\textsuperscript{6}，正思戎马泪盈巾\textsuperscript{7}。
\end{poem}

\begin{annotations}
\notetext{1}{扑枣：用杆子敲打枣树让枣子落下。杜甫住草堂时任凭西邻的老妇人前来打枣充饥。}
\notetext{2}{不为困穷宁有此：不是因为穷困到了极点，怎么会做这样的事？宁，怎么。}
\notetext{3}{只缘恐惧转须亲：正因为她内心恐惧自卑，反而更要亲近她。杜甫的同情不是居高临下，是设身处地。}
\notetext{4}{即防远客虽多事：她防备你这个新来的远客，虽然是她想多了。}
\notetext{5}{便插疏篱却甚真：但你立刻插上篱笆，却未免太当真了。杜甫委婉批评吴郎的行为伤害了老妇人。}
\notetext{6}{征求贫到骨：官府的横征暴敛让她穷到骨头里。征求，征敛索取。}
\notetext{7}{正思戎马泪盈巾：想到战乱还在继续，泪水湿透了衣巾。由一个人的苦难看到了天下人的苦难。}
\end{annotations}

\begin{appreciation}
这首诗写于大历二年（767年），杜甫在夔州瀼西草堂居住时。西邻有一位贫苦的老妇人，常来堂前打枣充饥。后来杜甫搬走了，把草堂借给亲戚吴郎住。吴郎一来就插上了篱笆，不让老妇人靠近。杜甫知道后，写了这首诗劝吴郎。

堂前扑枣任西邻，无食无儿一妇人——我住在堂前时，任凭西邻的老妇人来打枣；她没有食物，没有儿女，孤苦一人。扑枣就是打枣，用杆子敲打枣树让枣子落下。任字写出了杜甫的态度：不是默许，是主动允许；不是无奈，是心甘情愿。

不为困穷宁有此？只缘恐惧转须亲——不是因为穷困到极点，怎么会做这样的事？正因为她内心恐惧、自卑，我们反而更要亲近她。宁有此，是反问：如果不是走投无路，谁愿意来别人家打枣呢？转须亲，是杜甫最动人的地方：他不责备老妇人的唐突，反而觉得应该更亲近她，因为她已经够害怕了。

即防远客虽多事，便插疏篱却甚真——她防备你这个远客，虽然是多心；但你插上稀疏的篱笆，却未免太当真了。防远客，指老妇人对吴郎这个陌生人的戒备。多事，是说她想多了。但便插疏篱却甚真——你吴郎立刻就插篱笆，这反应未免太真实、太伤人了。杜甫的语气很委婉，但批评的意思很清楚。

已诉征求贫到骨，正思戎马泪盈巾——她已经告诉我，官府的征敛让她穷到了骨头里；想到战乱不止，我不禁泪湿衣巾。征求是官府的横征暴敛，贫到骨是穷到了极点。戎马指战乱，杜甫由一个人的苦难，想到了天下人的苦难。泪盈巾，不是为老妇人而哭，是为这个时代而哭。

这首诗最打动人的，是杜甫的推己及人。他不是在施舍，是在尊重；不是在怜悯，是在共情。一个无食无儿的妇人，在杜甫眼中不是麻烦，是需要被温柔以待的人。这种胸怀，穿越千年，依然让人动容。
\end{appreciation}
